% ===================================================================== PREAMBULE COMMUN — Carnet d'entraînement Fiche 9
\documentclass[11pt,a4paper]{article}
\usepackage[utf8]{inputenc}
\usepackage[T1]{fontenc}
\usepackage{lmodern}
\usepackage[french]{babel}
\usepackage{amsmath,amssymb,mathtools}
\usepackage{icomma}
\usepackage{siunitx}
\sisetup{output-decimal-marker={,},per-mode=symbol}
\DeclareSIUnit{\molar}{M}
\DeclareSIUnit{\Molar}{mol\,L^{-1}}
\usepackage[version=4]{mhchem}
\usepackage{xcolor}
\usepackage{colortbl}
\usepackage{tikz}
\usepackage{pgfplots}
\usepackage[most]{tcolorbox}
\usepackage{enumitem}
\usepackage{array}
\usepackage{booktabs}
\usepackage{multicol}
\usepackage{fancyhdr}
\usepackage[hidelinks]{hyperref}
\usetikzlibrary{arrows.meta,positioning,calc,decorations.pathmorphing,
  decorations.markings,angles,quotes,patterns,backgrounds,
  shapes.geometric,shapes.misc,fadings,intersections,fit}
\pgfplotsset{compat=1.18}
\usepackage[a4paper,margin=2.2cm,headheight=26pt,headsep=10pt]{geometry}

% ---------- Palette ----------
\definecolor{primary}{RGB}{150,98,162}
\definecolor{primarydark}{RGB}{92,52,110}
\definecolor{defcol}{RGB}{0,114,178}
\definecolor{thmcol}{RGB}{0,140,120}
\definecolor{formulacol}{RGB}{198,93,9}
\definecolor{remarkcol}{RGB}{120,94,200}
\definecolor{attncol}{RGB}{200,40,55}
\definecolor{excol}{RGB}{0,135,90}
\definecolor{methcol}{RGB}{90,90,100}
\definecolor{cationcol}{RGB}{0,90,190}
\definecolor{anioncol}{RGB}{200,60,55}
\definecolor{cristalcol}{RGB}{110,105,120}
\definecolor{eaucol}{RGB}{120,180,220}
\definecolor{solcol}{RGB}{0,135,90}
\definecolor{kpscol}{RGB}{198,93,9}
\definecolor{qcol}{RGB}{120,94,200}
\definecolor{precipcol}{RGB}{110,70,40}
\definecolor{exercol}{RGB}{150,98,162}
\definecolor{corrcol}{RGB}{0,120,100}

% ---------- Macros ----------
\newcommand{\Kps}{K_{\mathrm{ps}}}
\newcommand{\Ce}[1]{\left[\,\ce{#1}\,\right]}

% ---------- Style des boîtes ----------
\tcbset{boxbase/.style={enhanced,breakable,boxrule=0.9pt,arc=2.5pt,
   left=3mm,right=3mm,top=2.3mm,bottom=2.3mm,
   fonttitle=\bfseries,coltitle=white,toptitle=0.6mm,bottomtitle=0.6mm}}

\newtcolorbox[auto counter]{definition}[1][]{%
  boxbase,colback=defcol!5,colframe=defcol,
  title={\textsc{Définition}~\thetcbcounter\ \ifx\relax#1\relax\relax\else\textmd{--- \itshape #1}\fi}}
\newtcolorbox[auto counter]{theoreme}[1][]{%
  boxbase,colback=thmcol!6,colframe=thmcol,
  title={\textsc{Loi / Propriété}~\thetcbcounter\ \ifx\relax#1\relax\relax\else\textmd{--- \itshape #1}\fi}}
\newtcolorbox[auto counter]{formule}[1][]{%
  boxbase,colback=formulacol!6,colframe=formulacol,
  title={\textsc{Formule}~\thetcbcounter\ \ifx\relax#1\relax\relax\else\textmd{--- \itshape #1}\fi}}
\newtcolorbox{remarque}[1][]{%
  boxbase,colback=remarkcol!6,colframe=remarkcol,
  title={\textsc{Remarque}\ifx\relax#1\relax\relax\else\textmd{ : \itshape #1}\fi}}
\newtcolorbox{attention}[1][]{%
  boxbase,colback=attncol!6,colframe=attncol,
  title={\textsc{Attention}\ifx\relax#1\relax\relax\else\textmd{ : \itshape #1}\fi}}
\newtcolorbox{exemple}[1][]{%
  boxbase,colback=excol!5,colframe=excol,
  title={\textsc{Exemple}\ifx\relax#1\relax\relax\else\textmd{ : \itshape #1}\fi}}
\newtcolorbox{methode}[1][]{%
  boxbase,colback=methcol!7,colframe=methcol,sharp corners,
  title={\textsc{Méthode}\ifx\relax#1\relax\relax\else\textmd{ : \itshape #1}\fi}}
\newtcolorbox{astuce}[1][]{%
  boxbase,colback=primary!7,colframe=primary,
  title={\textsc{Astuce}\ifx\relax#1\relax\relax\else\textmd{ : \itshape #1}\fi}}

\newtcolorbox[auto counter]{exercice}[1][]{%
  boxbase,colback=exercol!4,colframe=exercol,
  title={\textsc{Exercice}~\thetcbcounter\ \ifx\relax#1\relax\relax\else\textmd{--- \itshape #1}\fi}}
\newtcolorbox[auto counter]{correction}[1][]{%
  boxbase,colback=corrcol!4,colframe=corrcol,
  title={\textsc{Corrigé}~\thetcbcounter\ \ifx\relax#1\relax\relax\else\textmd{--- \itshape #1}\fi}}

% ---------- Environnement « où : » ----------
\newenvironment{ou}{%
  \par\smallskip\noindent\textbf{où :}\nopagebreak
  \begin{itemize}[leftmargin=1.8em,topsep=2pt,itemsep=1.5pt,parsep=0pt]}%
  {\end{itemize}}

% ---------- Styles TikZ ----------
\tikzset{
  cat/.style={circle,fill=cationcol,draw=cationcol!50!black,inner sep=0pt,
              minimum size=5mm,font=\bfseries\scriptsize,text=white},
  an/.style={circle,fill=anioncol,draw=anioncol!50!black,inner sep=0pt,
             minimum size=5mm,font=\bfseries\scriptsize,text=white},
  crx/.style={circle,fill=cristalcol,draw=black!50,inner sep=0pt,
              minimum size=5mm,font=\bfseries\scriptsize,text=white},
  ptn/.style={circle,fill=black,inner sep=1.1pt}}

% ---------- En-têtes ----------
\pagestyle{fancy}
\fancyhf{}
\fancyhead[L]{\small\textcolor{primarydark}{\textbf{Fiche 9} — Solubilité et produit de solubilité}}
\fancyhead[R]{\small\textcolor{primarydark}{\textsc{Carnet d'entraînement}}}
\fancyfoot[C]{\small\thepage}
\renewcommand{\headrulewidth}{0.8pt}
\renewcommand{\headrule}{\hbox to\headwidth{\color{primary}\leaders\hrule height \headrulewidth\hfill}}
\usepackage{titlesec}
\titleformat{\section}{\Large\bfseries\color{primarydark}}{\thesection}{0.6em}{}
\titleformat{\subsection}{\large\bfseries\color{primary}}{\thesubsection}{0.6em}{}

% ---------- Bandeau de titre ----------
% #1 = thème/étiquette   #2 = titre principal   #3 = sous-titre
\newcommand{\bandeau}[3]{%
\begin{center}
\begin{tikzpicture}
\node[rectangle,rounded corners=7pt,fill=primary,text=white,
      minimum width=15.6cm,minimum height=1.95cm,align=center,inner sep=8pt]
  {{\large\bfseries #1}\\[3pt]{\Large\bfseries #2}\\[3pt]{\small #3}};
\end{tikzpicture}
\end{center}
\vspace{2mm}}
\begin{document}
\bandeau{Thème 1 — Équation de dissolution \& produit de solubilité}%
{Tuto — la mécanique de base}%
{Écrire l'équation de dissolution, puis l'expression littérale du $\Kps$}

\noindent Tout problème de solubilité commence par \textbf{deux gestes purement mécaniques} (aucun
calcul, aucune donnée numérique) : \textbf{(1)} écrire l'\emph{équation de dissolution} du sel, puis
\textbf{(2)} en déduire l'\emph{expression du $\Kps$}. Les maîtriser, c'est faire « la moitié du
travail » de tout exercice de la fiche.

\section*{1. L'équation de dissolution}
Un composé ionique solide plongé dans l'eau se sépare en ses ions hydratés. Pour un sel
$\mathrm{M}_p\mathrm{X}_q$ (cation $\mathrm{M}^{q+}$, anion $\mathrm{X}^{p-}$) :
\[ \boxed{\ \ce{M_{p}X_{q}(s) <=> p\,M^{q+}(aq) + q\,X^{p-}(aq)}\ } \]

\begin{methode}[3 réflexes]
\textbf{(i)} solide $(s)$ à gauche, double flèche $\ce{<=>}$ ; \textbf{(ii)} ions $(aq)$ à droite, avec
leur charge ; \textbf{(iii)} reporter les \emph{indices} de la formule en \emph{coefficients}
($p$ cations, $q$ anions), puis \textbf{vérifier la neutralité} des charges.
\end{methode}

\begin{attention}[l'erreur n\textsuperscript{o}\,1 : oublier les coefficients]
\ce{Ag3PO4} libère \textbf{3} ions \ce{Ag+} et \textbf{1} ion \ce{PO4^{3-}}. Un groupement polyatomique
(\ce{PO4^{3-}}, \ce{CO3^{2-}}, \ce{SO4^{2-}}, \ce{OH-}\dots) reste \textbf{entier} : il ne se casse pas.
Le « 3 » deviendra un \textbf{exposant} dans le $\Kps$ — d'où son importance.
\end{attention}

\section*{2. L'expression du produit de solubilité $\Kps$}
Le $\Kps$ est la \textbf{constante d'équilibre} de la dissolution : produit des concentrations des ions,
chacune \textbf{élevée à son coefficient}.
\[ \boxed{\ \Kps = [\mathrm{M}^{q+}]^{\,p}\,[\mathrm{X}^{p-}]^{\,q}\ } \]

\begin{astuce}[la règle d'or]
\emph{L'exposant de chaque ion dans le $\Kps$ = son coefficient dans l'équation.} Le \textbf{solide
n'apparaît jamais} (activité $=1$ : que l'on ait \qty{1}{\milli\gram} ou \qty{1}{\kilo\gram} de solide,
sa contribution est la même).
\end{astuce}

\begin{exemple}[les deux gestes, sel par sel]
\begin{center}\renewcommand{\arraystretch}{1.35}
\begin{tabular}{l l c l}
\toprule
\rowcolor{primary!12}
\textbf{Sel} & \textbf{Équation de dissolution} & \textbf{Type} & \textbf{$\Kps$}\\
\midrule
\ce{AgCl}     & $\ce{AgCl(s) <=> Ag+ + Cl-}$              & $1{:}1$ & $\Ce{Ag+}\Ce{Cl-}$\\
\ce{CaF2}     & $\ce{CaF2(s) <=> Ca^{2+} + 2 F-}$        & $1{:}2$ & $\Ce{Ca^2+}\Ce{F-}^{2}$\\
\ce{Ag2CrO4}  & $\ce{Ag2CrO4(s) <=> 2 Ag+ + CrO4^{2-}}$ & $2{:}1$ & $\Ce{Ag+}^{2}\Ce{CrO4^2-}$\\
\ce{Fe(OH)3}  & $\ce{Fe(OH)3(s) <=> Fe^{3+} + 3 OH-}$   & $1{:}3$ & $\Ce{Fe^3+}\Ce{OH-}^{3}$\\
\ce{Ag3PO4}   & $\ce{Ag3PO4(s) <=> 3 Ag+ + PO4^{3-}}$   & $3{:}1$ & $\Ce{Ag+}^{3}\Ce{PO4^3-}$\\
\ce{Ca3(PO4)2}& $\ce{Ca3(PO4)2(s) <=> 3 Ca^{2+} + 2 PO4^{3-}}$ & $3{:}2$ & $\Ce{Ca^2+}^{3}\Ce{PO4^3-}^{2}$\\
\bottomrule
\end{tabular}
\end{center}
Le \textbf{type} $p{:}q$ (cations $:$ anions) résume tout : il fixe la forme du $\Kps$ et, plus tard, la
relation avec la solubilité (Thème 2).
\end{exemple}

\begin{attention}[le $\Kps$ est \emph{sans unité}, mais dépend de $T$]
Les concentrations sont en \si{\mole\per\litre}, mais le $\Kps$ (constante d'équilibre) s'écrit
\textbf{sans unité} : $\Kps(\ce{AgCl})=\num{1{,}8e-10}$. En revanche il \textbf{dépend de la
température} : la dissolution étant souvent endothermique, $\Kps$ croît en général avec $T$.
\end{attention}

\begin{remarque}[cas particulier $MX_n$ (formule du livre)]
Pour un sel $\mathrm{MX}_n$ (un cation, $n$ anions) : $\ce{MX_{n}(s) <=> M^{n+} + n X-}$ et
$\Kps=[\mathrm{M}^{n+}][\mathrm{X}^-]^{n}$. La forme générale
$\Kps=[\mathrm{M}^{q+}]^{p}[\mathrm{X}^{p-}]^{q}$ couvre en plus \ce{Ag2CrO4}, \ce{Ca3(PO4)2}, etc.
\medskip

\noindent\emph{Objectif : lire un nom ou une formule $\to$ écrire l'équation $\to$ écrire le $\Kps$,
sans hésitation et sans calcul.}
\end{remarque}

\end{document}
